\documentclass[12pt,a4paper]{report}
\usepackage[spanish,es-nodecimaldot]{babel}
\usepackage[utf8]{inputenc}
\usepackage[T1]{fontenc}
\usepackage{lmodern}
\usepackage{graphicx}
\usepackage{hyperref}
\usepackage{geometry}
\usepackage{setspace}
\usepackage{float}
\usepackage{longtable}
\usepackage{booktabs}
\usepackage{tabularx}
\usepackage{array}
\usepackage{enumitem}
\usepackage{tikz}
\usetikzlibrary{positioning,arrows.meta}
\usepackage{xcolor}
\usepackage{listings}
\usepackage{fancyhdr}
\usepackage{microtype}

\geometry{margin=2.5cm}
\setlength{\headheight}{15pt}
\onehalfspacing
\setlength{\parindent}{1.25cm}
\setlength{\parskip}{0.35em}
\hypersetup{
  colorlinks=true,
  linkcolor=black,
  urlcolor=blue,
  citecolor=black,
  pdftitle={Pipeline CI/CD Automatizado para una API REST},
  pdfauthor={Francisco Giovanni López Moscosa}
}
\pagestyle{fancy}
\fancyhf{}
\fancyhead[L]{Elementos de DevOps}
\fancyhead[R]{Francisco Giovanni López Moscosa}
\fancyfoot[C]{\thepage}
\lstset{basicstyle=\ttfamily\small,breaklines=true,frame=single,columns=fullflexible}
\newcommand{\conceptbox}[2]{\begin{center}\fbox{\begin{minipage}{0.92\textwidth}\textbf{#1.} #2\end{minipage}}\end{center}}
\definecolor{github}{HTML}{24292F}
\definecolor{docker}{HTML}{2496ED}
\definecolor{aws}{HTML}{FF9900}
\definecolor{success}{HTML}{2DA44E}

\begin{document}
\spanishdeactivate{>}

%------------------------------------------------
% PORTADA INSTITUCIONAL
\begin{titlepage}
\noindent
\begin{minipage}{0.45\textwidth}
\raggedright
\IfFileExists{images.jpg}{\includegraphics[width=0.3\textwidth]{images.jpg}}{\fbox{\parbox[c][2.5cm][c]{3cm}{\centering Logo\\institucional}}}
\end{minipage}
\vspace{1.5cm}
\centering
{\scshape\LARGE Universidad Tecnológica de Querétaro \par}
\vspace{2cm}
{\huge\bfseries Elementos de DevOps\par}
\vspace{0.8cm}
{\LARGE\bfseries Proyecto Integrador: Pipeline CI/CD Automatizado para una API REST\par}
\vspace{2cm}
{\Large Alumno: Francisco Giovanni López Moscosa \par}
{\Large Carrera: Ingeniería en Desarrollo de Software Multiplataforma \par}
{\Large Matrícula: 2023371156 \par}
\vfill
{\Large Docente: Emmanuel Martínez Hernández \par}
\vspace{0.5cm}
{\Large Querétaro, Qro. \today \par}
\end{titlepage}

\tableofcontents
\listoffigures
\listoftables
\newpage

\chapter{Introducción}

El desarrollo de una aplicación no termina cuando el código compila en la computadora del programador. Una API REST que será utilizada por clientes, aplicaciones web o servicios externos debe pasar por controles de calidad repetibles, empaquetarse con todas sus dependencias y llegar al servidor de producción de una forma trazable. Cuando estas actividades se realizan manualmente, cada entrega depende de instrucciones informales, aumenta el riesgo de omitir una prueba y se vuelve difícil conocer qué versión exacta se encuentra en ejecución. Las prácticas de integración continua y entrega continua, conocidas como CI/CD por sus siglas en inglés, atienden este problema al definir una secuencia automática y verificable desde un cambio versionado hasta una versión disponible para los usuarios \cite{githubactions}.

El presente proyecto integra una API REST de gestión de tareas y categorías con Docker, GitHub Actions, Docker Hub y Amazon EC2. La aplicación fue desarrollada con ASP.NET Core 10 y emplea SQLite como almacenamiento local. Sus endpoints permiten consultar y administrar tareas y categorías, descargar un respaldo de la base de datos, vaciar los datos y conocer el estado del servicio mediante \texttt{/api/health}. Esta última ruta se utiliza como una señal objetiva de disponibilidad durante el despliegue. El proyecto incluye pruebas de integración escritas con Jest, las cuales inician la API en un puerto temporal y usan una base SQLite independiente para no alterar la información de desarrollo.

La calidad se controla antes de publicar una imagen. En cada cambio dirigido a la rama \texttt{main}, GitHub Actions restaura dependencias de .NET y Node.js, ejecuta las pruebas y mide cobertura sobre el proceso de ASP.NET Core con \texttt{dotnet-coverage}. La política del proyecto exige al menos 70\% de cobertura de líneas. El resultado observado fue 19 pruebas aprobadas y 96.92\% de cobertura, por lo que el cambio fue apto para continuar al empaquetado. Esta condición evita que una imagen sin validación automática avance al entorno público y ofrece evidencia reproducible de la calidad de cada entrega.

Después de las pruebas, Docker construye la aplicación mediante un archivo multi-stage. La primera etapa utiliza el SDK para restaurar, compilar y publicar; la segunda conserva únicamente el runtime necesario para ejecutar el ensamblado. La imagen se registra en Docker Hub con dos etiquetas: \texttt{latest}, útil para identificar la versión más reciente, y el SHA completo del commit, que representa una versión inmutable. El identificador de commit permite relacionar la imagen, la ejecución del workflow y el código fuente que se desplegó. Por ello, ante una incidencia es posible identificar y recuperar una versión específica sin depender de una etiqueta mutable \cite{docker}.

La entrega continua se realiza en una instancia Ubuntu de Amazon EC2. GitHub Actions se autentica por SSH usando una clave privada almacenada exclusivamente como Secret y verifica la clave pública del host antes de conectarse. En el servidor, un script descarga la imagen identificada por SHA, la inicia en el puerto local alterno, consulta el endpoint de salud y redirige Nginx sólo cuando la respuesta es satisfactoria. Esta estrategia blue/green conserva el contenedor anterior mientras se valida el nuevo y reduce la interrupción del servicio. Los datos SQLite se alojan en un volumen Docker separado de los contenedores, por lo que se conservan aunque cambie la imagen de la aplicación. El enfoque combina automatización, seguridad de secretos y disponibilidad para demostrar un ciclo de vida profesional de una API REST \cite{aws}.

\chapter{Objetivos}

\section{Objetivo general}
Implementar un pipeline CI/CD para una API REST que automatice sus pruebas, cobertura, contenedorización, publicación de imagen Docker y despliegue seguro en una instancia AWS EC2.

\section{Objetivos específicos}
\begin{enumerate}[label=\arabic*.]
  \item Ejecutar pruebas de integración de la API de tareas y categorías en cada cambio enviado a la rama \texttt{main}.
  \item Generar y validar un reporte de cobertura de código con un umbral mínimo de 70\%.
  \item Empaquetar la aplicación en una imagen Docker multi-stage y excluir archivos innecesarios del contexto de construcción.
  \item Publicar la imagen en Docker Hub con las etiquetas \texttt{latest} y el SHA del commit.
  \item Desplegar la imagen validada en Amazon EC2 mediante SSH y exponer la API con Nginx en el puerto HTTP 80.
  \item Preservar los datos de SQLite y reducir la interrupción usando contenedores blue/green en puertos locales alternos.
\end{enumerate}

\chapter{Arquitectura y desarrollo}

\section{Componentes de la solución}

La Tabla~\ref{tab:components} resume los componentes utilizados y su responsabilidad dentro de la solución. La separación de responsabilidades permite que el servicio de integración continua se encargue de la validación, Docker Hub funcione como registro de artefactos y EC2 ejecute únicamente imágenes ya aprobadas por el pipeline.

\begin{table}[H]
\centering
\caption{Componentes de la arquitectura.}
\label{tab:components}
\begin{tabularx}{\textwidth}{>{\raggedright\arraybackslash}p{3.2cm}X}
\toprule
\textbf{Componente} & \textbf{Responsabilidad} \\
\midrule
ASP.NET Core 10 y SQLite & Implementar la API REST y persistir tareas y categorías. \\
Jest y dotnet-coverage & Ejecutar pruebas de integración y obtener cobertura sobre el proceso backend. \\
GitHub Actions & Orquestar CI, publicar la imagen y ejecutar el despliegue autorizado. \\
Docker Hub & Almacenar imágenes identificadas por \texttt{latest} y SHA. \\
AWS EC2 Ubuntu & Ejecutar Docker, Nginx y la aplicación en producción. \\
Nginx & Recibir HTTP en el puerto 80 y dirigir el tráfico al contenedor activo. \\
\bottomrule
\end{tabularx}
\end{table}

\section{Flujo CI/CD}

La Figura~\ref{fig:architecture} presenta la arquitectura implementada. Un \texttt{push} o \texttt{pull request} dirigido a \texttt{main} inicia el trabajo de pruebas. Sólo los \texttt{push} exitosos a la rama principal pueden usar los Secrets de producción para publicar y desplegar. De este modo, una solicitud de cambios no expone token de Docker Hub, llave SSH ni dirección del servidor.

\begin{figure}[H]
\centering
\begin{tikzpicture}[
  node distance=8mm and 9mm,
  box/.style={draw, rounded corners, minimum width=2.55cm, minimum height=1.1cm, align=center, font=\small},
  arrow/.style={-{Latex[length=2mm]}, thick}
]
\node[box, fill=github!12] (git) {GitHub\\\texttt{push} a \texttt{main}};
\node[box, fill=success!12, right=of git] (ci) {GitHub Actions\\Pruebas y cobertura};
\node[box, fill=docker!12, right=of ci] (hub) {Docker Hub\\\texttt{latest} y SHA};
\node[box, fill=aws!18, right=of hub] (ec2) {AWS EC2\\Docker y Nginx};
\node[box, fill=success!12, below=of ec2] (api) {API pública\\\texttt{/api/health}};
\draw[arrow] (git) -- (ci);
\draw[arrow] (ci) -- node[above, font=\scriptsize]{si cobertura $\geq$ 70\%} (hub);
\draw[arrow] (hub) -- node[above, font=\scriptsize]{SSH} (ec2);
\draw[arrow] (ec2) -- (api);
\end{tikzpicture}
\caption{Arquitectura del pipeline CI/CD implementado.}
\label{fig:architecture}
\end{figure}

\section{Pruebas y cobertura}

Las pruebas de integración no utilizan la base de datos de desarrollo. Antes de iniciar la API, Jest elimina la base temporal \texttt{webapp.test.db}; durante la suite crea categorías y tareas, consulta recursos existentes e inexistentes, actualiza y elimina registros, descarga un respaldo y vacía la información. Los casos incluyen respuestas exitosas y respuestas de error \texttt{400} y \texttt{404}. El endpoint \texttt{/api/health} también se prueba para comprobar que el servicio está disponible.

El workflow ejecuta \texttt{dotnet-coverage collect "npm test"} y publica el archivo Cobertura como artefacto. Un paso posterior lee la tasa de líneas y falla si no alcanza 70\%. Por tanto, la cobertura no es un dato informativo opcional, sino una condición para continuar a la publicación de la imagen.

\section{Contenedorización y despliegue}

El \texttt{Dockerfile} utiliza compilación multi-stage. Los archivos de compilación, pruebas, repositorio Git, bases locales, variables de entorno y reportes se excluyen mediante \texttt{.dockerignore}. La imagen final expone el puerto 80 y ejecuta únicamente \texttt{WebApp.dll}. En producción, el volumen \texttt{webapp-api-data} se monta en \texttt{/data}; la variable \texttt{DatabasePath} dirige el archivo SQLite a ese volumen persistente.

El despliegue identifica cuál puerto local está activo entre 8081 y 8082. A continuación inicia la versión candidata en el puerto opuesto, realiza hasta 30 intentos contra \texttt{/api/health}, genera la configuración de Nginx y recarga el proxy. Sólo después elimina el contenedor anterior. La Figura~\ref{fig:bluegreen} ilustra este procedimiento.

\begin{figure}[H]
\centering
\begin{tikzpicture}[
  node distance=9mm and 14mm,
  box/.style={draw, rounded corners, minimum width=3.2cm, minimum height=0.95cm, align=center, font=\small},
  arrow/.style={-{Latex[length=2mm]}, thick}
]
\node[box, fill=aws!18] (nginx) {Nginx :80};
\node[box, fill=docker!12, right=of nginx] (active) {Blue :8081\\versión activa};
\node[box, fill=success!12, below=of active] (candidate) {Green :8082\\versión candidata};
\node[box, fill=github!12, left=of candidate] (health) {Consulta\\\texttt{/api/health}};
\draw[arrow] (nginx) -- (active);
\draw[arrow] (candidate) -- (health);
\draw[arrow, dashed] (health) -| node[pos=0.25, below, font=\scriptsize]{respuesta 200} (nginx);
\end{tikzpicture}
\caption{Despliegue blue/green con validación de salud antes de redirigir Nginx.}
\label{fig:bluegreen}
\end{figure}

\chapter{Resultados}

\section{Ejecución de calidad}

La primera etapa del pipeline se ejecutó correctamente. La suite reportó 19 pruebas aprobadas, sin pruebas fallidas ni omitidas. La compilación Release de ASP.NET Core finalizó sin errores ni advertencias. La cobertura de líneas obtenida sobre el backend fue de 96.92\%, valor que supera en 26.92 puntos porcentuales el umbral exigido de 70\%. Este resultado demuestra que la suite cubre las operaciones principales de la API y sus escenarios de error.

La Figura~\ref{fig:quality-results} resume los indicadores observados. La barra verde representa el valor obtenido y la marca oscura indica el mínimo configurado en el workflow. Debido a que el valor obtenido es superior al requisito, GitHub Actions habilitó el trabajo de publicación y despliegue.

\begin{figure}[H]
\centering
\begin{tikzpicture}[x=0.10cm,y=0.75cm]
\draw[->] (0,0) -- (105,0) node[right]{Cobertura de líneas (\%)};
\foreach \x in {0,10,...,100} \draw (\x,0.12) -- (\x,-0.12) node[below, font=\scriptsize]{\x};
\fill[success!65] (0,0.45) rectangle (96.92,1.35);
\draw[thick] (0,0.45) rectangle (96.92,1.35);
\draw[very thick, github] (70,0.2) -- (70,1.6);
\node[above, font=\small] at (96.92,1.35) {96.92\% obtenido};
\node[below, font=\scriptsize, github] at (70,0.2) {mínimo 70\%};
\end{tikzpicture}
\caption{Cobertura de líneas obtenida en la ejecución de CI.}
\label{fig:quality-results}
\end{figure}

\section{Publicación de la imagen}

Al completar las pruebas, Docker Hub recibió la imagen \texttt{fnxi/webapp-api}. El proceso publicó la etiqueta \texttt{latest} y una etiqueta inmutable con el SHA \texttt{4f2ed224fc739b84b8eb0d26e5f0b10f6e38c846}. Esta segunda etiqueta relaciona directamente el código fuente con el artefacto desplegado. La publicación se realizó mediante \texttt{docker/login-action} y el token se obtuvo de GitHub Secrets, por lo que ningún valor confidencial fue incluido en el repositorio.

\section{Despliegue y disponibilidad}

La etapa de despliegue se conectó a la EC2 por SSH, descargó la imagen identificada por SHA y creó el contenedor \texttt{webapp-api-blue}. La inspección de la instancia confirmó que el contenedor está activo, que ejecuta la misma imagen publicada y que se encuentra mapeado como \texttt{127.0.0.1:8081->80/tcp}. Nginx recibe las solicitudes externas en el puerto 80 y las dirige al puerto local activo.

Finalmente, se consultó la URL pública \url{http://3.14.246.106/api/health}. La respuesta recibida fue \texttt{HTTP 200} con el cuerpo mostrado en el Listado~\ref{lst:health}. Este resultado valida que la ruta HTTP, Nginx, Docker, la API y la instancia EC2 están integrados correctamente.

\begin{lstlisting}[caption={Respuesta del endpoint público de salud.},label={lst:health}]
{"statusCode":200,"data":{"status":"healthy"}}
\end{lstlisting}

\begin{figure}[H]
\centering
\begin{tikzpicture}[
  node distance=8mm and 10mm,
  box/.style={draw, rounded corners, minimum width=3.4cm, minimum height=1cm, align=center, font=\small},
  arrow/.style={-{Latex[length=2mm]}, thick}
]
\node[box, fill=docker!12] (image) {Imagen Docker Hub\\SHA \texttt{4f2ed224...}};
\node[box, fill=success!12, right=of image] (container) {\texttt{webapp-api-blue}\\\texttt{127.0.0.1:8081->80}};
\node[box, fill=aws!18, right=of container] (public) {EC2 pública\\HTTP 200 en \texttt{/api/health}};
\draw[arrow] (image) -- (container);
\draw[arrow] (container) -- (public);
\end{tikzpicture}
\caption{Evidencia del artefacto publicado, contenedor desplegado y disponibilidad pública.}
\label{fig:deployment-result}
\end{figure}

\section{Seguridad y trazabilidad}

La información sensible se gestionó con GitHub Secrets: usuario y token de Docker Hub, host y usuario SSH de EC2, llave privada SSH y clave pública del servidor. El workflow activa la publicación y el despliegue sólo ante un \texttt{push} a \texttt{main}; las solicitudes de cambio ejecutan validaciones sin acceder a secretos. Además, la conexión SSH usa \texttt{StrictHostKeyChecking=yes}, por lo que no acepta un servidor cuya clave no coincida con el valor almacenado. Estas decisiones reducen la exposición de credenciales y mitigan el riesgo de desplegar desde una fuente no verificada \cite{githubsecrets}.

\chapter{Conclusiones}

El proyecto demostró que es posible automatizar el ciclo completo de entrega de una API REST con herramientas accesibles y ampliamente utilizadas. GitHub Actions ejecuta pruebas de integración y aplica un criterio de cobertura de 70\% antes de publicar una imagen. El resultado de 19 pruebas exitosas y 96.92\% de cobertura aporta evidencia objetiva de que la versión construida satisface los controles definidos. La imagen multi-stage reduce el contenido de producción y las etiquetas \texttt{latest} y SHA permiten distinguir la versión reciente de una versión exacta e inmutable.

El despliegue en EC2 validó la integración de Docker, Nginx, SSH y GitHub Secrets. El contenedor \texttt{webapp-api-blue} se inició desde la imagen publicada y la respuesta HTTP 200 de \texttt{/api/health} confirmó que la API está disponible en la dirección pública. El patrón blue/green permite iniciar y probar un contenedor candidato antes de redirigir el tráfico, lo que reduce el riesgo de indisponibilidad frente a un reemplazo directo. Asimismo, el volumen Docker mantiene la base SQLite separada del ciclo de vida del contenedor.

Como mejora futura, el dominio de la API puede ampliarse hasta cubrir el número de endpoints solicitado por la rúbrica, incorporando recursos adicionales y sus pruebas. También sería conveniente migrar SQLite a un servicio de base de datos administrado si aumenta la concurrencia, usar HTTPS con un certificado TLS y sustituir la exposición global de SSH por un runner propio o AWS Systems Manager. Estas mejoras conservarían el flujo CI/CD implementado mientras fortalecen escalabilidad, seguridad y operación continua.

\begin{thebibliography}{9}
\bibitem{githubactions}
GitHub. \textit{Understanding GitHub Actions}. Consultado el 7 de octubre de 2026. \url{https://docs.github.com/actions/about-github-actions/understanding-github-actions}.

\bibitem{docker}
Docker. \textit{Docker documentation: Build images and registry}. Consultado el 7 de octubre de 2026. \url{https://docs.docker.com/}.

\bibitem{aws}
Amazon Web Services. \textit{Amazon EC2 User Guide for Linux Instances}. Consultado el 7 de octubre de 2026. \url{https://docs.aws.amazon.com/AWSEC2/latest/UserGuide/}.

\bibitem{githubsecrets}
GitHub. \textit{Using secrets in GitHub Actions}. Consultado el 7 de octubre de 2026. \url{https://docs.github.com/actions/security-for-github-actions/security-guides/using-secrets-in-github-actions}.

\bibitem{aspnet}
Microsoft. \textit{ASP.NET Core documentation}. Consultado el 7 de octubre de 2026. \url{https://learn.microsoft.com/aspnet/core/}.
\end{thebibliography}

\end{document}
